\documentclass[10pt,a4paper]{amsart}
\usepackage[margin=19mm]{geometry}
\usepackage{amsmath,amssymb,mathtools}
\usepackage[scr=rsfs,cal=euler]{mathalfa}
\usepackage{xfrac}
\usepackage[unicode,hidelinks,bookmarks=false]{hyperref}
\newcommand{\ie}{i.e.\ }
\newcommand{\cf}{cf.\ }
\newcommand{\resp}{respectively\ }
\newcommand{\Subsection}{\subsection}
\providecommand{\nobreakdash}{\nobreak\hbox{-}}
\setlength{\emergencystretch}{2em}
\multlinegap0pt
\usepackage{mathtools}
\usepackage{amssymb}
\usepackage{bm}
\usepackage{algorithm}\usepackage{algpseudocode}
\usepackage[shortlabels]{enumitem}
\usepackage{tcolorbox}
\usepackage{tikz}
\usepackage{xcolor}
\newtheorem{proposition}{Proposition}
\newcommand{\bbR}{\mathbb{R}}
\newcommand{\bfx}{{\bf x}}
\newcommand{\bfy}{{\bf y}}
\newcommand{\nx}{{n_{x}}}
\newcommand{\ny}{{n_{y}}}
\newcommand{\pjoint}{\pi_{\rm joint}}
\newcommand{\pmarg}{\pi_{\rm marg}}
\title{Exploiting Exact Conditionals Improves Conditioning: \\ Provably Fast Mixing Time Bounds By Sampling from the Marginal}
\author{Abhijit Chowdhary, Federica Milinanni, Julianne Chung, Elizabeth Newman}
\date{Source: arXiv:2608.27884v1 (CC BY 4.0)}
\begin{document}
\maketitle

\noindent\emph{Source and licence.} Exploiting Exact Conditionals Improves Conditioning: \\ Provably Fast Mixing Time Bounds By Sampling from the Marginal. Version arXiv:2608.27884v1 (CC BY 4.0), distributed by arXiv under the Creative Commons Attribution 4.0 International licence, http://creativecommons.org/licenses/by/4.0/, as recorded in the arXiv licence metadata for this exact version and shown on the abstract page https://arxiv.org/abs/2608.27884v1. Full licence notice: \texttt{licenses/arxiv-2608.27884-CC-BY-4.0.txt}.\par\smallskip

\noindent\textit{Editorial scope.} Counted unit: the article's proposition prop:marginal-isoperimetry (source0.tex lines 1704--1713). The article's complete original proof of it is delivered with it (source0.tex lines 1715--1744, one \texttt{proof} environment of 30 lines). No mathematics is written, repaired or re-derived: the statement and the proof are the article's own lines, in the article's own notation, and no retained line is substituted or re-flowed.  Retained uncounted as the source-local context the proof needs: lines 1632--1635.  Nothing was omitted from any retained span, so the package carries no omission record.  No retained line contains a citation, so the wrapper carries no bibliography and no bibliography entry.  No retained line contains a figure environment, an image-inclusion command or drawing code, so no image is delivered and the package declares no asset dependency.  Editorial formatting only, in addition: an AMS \texttt{amsart} preamble that restates the article's own declarations and macros used by the retained lines, namely the environments \texttt{proposition} on separate counters, together with the standard packages those lines need.  The retained environments print in this document as Proposition 1 in order of appearance, whereas the article prints the same items with the numbering of its own full document.  The complete native source of the article as distributed in the arXiv e-print for this exact version is bundled unchanged as \texttt{sources/208746/source0.tex}.
  \begin{equation}
    \label{eq:isoperimetric-inequality}
    \pi(S_3) \geq \psi\,D(S_1,S_2)\,\pi(S_1)\pi(S_2).
  \end{equation}

\begin{proposition}[Marginalization Isoperimetric Bounds]
  \label{prop:marginal-isoperimetry}
  %
  If $\psi(\pjoint) < \infty$, then
  \begin{equation}
    \label{eq:marginal-isoperimetry}
    %
    \psi(\pjoint) \leq \psi(\pmarg).
  \end{equation}
\end{proposition}

\begin{proof}
  Let $(S_1,S_2,S_3)$ be a measurable partition of $\bbR^\ny$, and define the corresponding partition of the joint space by $A_i \coloneqq \bbR^\nx \times S_i$, for $i \in \{1,2,3\}$.
  %
  %
  %
  By the definition of the marginal, $\pjoint(A_i)=\pmarg(S_i)$ for each $i$.
  Moreover,
  \begin{equation*}
    D(A_1,A_2)
    = \inf_{\substack{\bfx_1,\bfx_2\in\bbR^\nx\\ \bfy_1\in S_1,\,\bfy_2\in S_2}}
    \left( \|\bfx_1-\bfx_2\|_2^2 + \|\bfy_1-\bfy_2\|_2^2 \right)^{1 / 2}
    = D(S_1,S_2),
  \end{equation*}
  where %
  the second equality is obtained by taking $\bfx_1=\bfx_2$.
  Thus, if $\psi$ is any admissible isoperimetric constant for $\pjoint$, applying~\eqref{eq:isoperimetric-inequality} to $(A_1,A_2,A_3)$ gives
  %
  %
  %
  %
  %
  \begin{equation*}
    \pmarg(S_3)  = \pjoint(A_3)
    \geq \psi\,D(A_1,A_2)\,\pjoint(A_1)\pjoint(A_2)
    = \psi\,D(S_1,S_2)\,\pmarg(S_1)\pmarg(S_2).
  \end{equation*}
  Hence, $\psi$ is also an admissible isoperimetric constant for $\pmarg$, and
  \eqref{eq:marginal-isoperimetry} follows by taking the supremum in $\psi$.
  %
\end{proof}

\end{document}
